\chapter{Prices as Information}\label{ch:b2:prices-as-information}

At an American racetrack, the odds on each horse are set by the bets themselves: the more money bet on
a horse, the less it pays when it wins. The odds are a price, and the price is a forecast. If the
forecast were right on average, a dollar bet on any kind of horse would return the same amount, the
dollar minus the track's share of the pool. The chart shows what a dollar bet to win actually returned
on all American races of 1992--2001.

\begin{openingchart}{Average rate of return per dollar bet to win, in per cent, on US horse races,
  1992--2001, by type of bet.}{Snowberg and Wolfers (2010)}
\begin{tikzpicture}
\begin{axis}[omrecord, xbar, bar width=13pt, xmin=-70, xmax=0, ymin=0.4, ymax=4.6, y dir=reverse,
    ytick=data, yticklabels from table={figdata/bachelor-2/08-prices-as-information/out/returns.csv}{bet},
    xlabel={Rate of return per dollar bet (per cent)}, grid=none, xmajorgrids,
    width=0.62\linewidth, height=0.42\linewidth, nodes near coords,
    every node near coord/.append style={font=\footnotesize, anchor=east}]
  \addplot[fill=omDef!70, draw=omDef] table[x=return, y=pos]
    {figdata/bachelor-2/08-prices-as-information/out/returns.csv};
\end{axis}
\end{tikzpicture}
\end{openingchart}

% figdata: returns_summary (gap 55.5, payback_favourite 0.945, payback_longest 0.39)
% ledger: b2:sw2010-returns
A dollar bet on the favourite in every race returned 94.5 cents on average; a dollar bet on a horse at
odds of 100/1 or longer returned 39 cents. Between the two ends lies a gap of 55.5 points. The prices
are good forecasts in one sense: favourites win far more often than longshots, as their odds say. They
are biased in another: measured against their odds, longshots win much less often than favourites do, so
the betting overvalues longshots and undervalues favourites. The
authors find that the pattern fits bettors who overweight small probabilities
(\cref{ch:b2:behavioural-evidence}) better than bettors who love risk.

\Cref{ch:b2:exchange-economies,ch:b2:welfare-theorems} treated prices as the terms on which goods
trade. This chapter treats them as messages. It counts what a planner would need to know to allocate
resources, shows how little a market needs instead, and asks how far a price can carry information
that no single trader has: completely, partly, or with a bias like the racetrack's. The answer depends
on what information costs, and it can be tested on markets whose prices are forecasts.

\begin{recall}
From the school volumes: prices signal scarcity and coordinate buyers and sellers. From the Year 1
volume: supply and demand at work, including price ceilings and the shortages they create. From
\cref{ch:b2:exchange-economies,ch:b2:welfare-theorems}: \omterm{def:b2:exchange-economies:competitive-equilibrium}{competitive equilibrium}, Pareto efficiency and
\omterm{def:b2:welfare-theorems:lump-sum}{lump-sum transfers}. From \cref{ch:b2:asset-prices}: asset prices and their unpredictability.
\end{recall}

\begin{notation}[Symbols of this chapter]\label{not:b2:prices-as-information:symbols}
$v$ is the value of an asset or the size of a harvest, unknown when trade takes place; $s_1, s_2$ and
$\theta$ are signals about it, and $\varepsilon$ is the part no signal predicts. $a$ is the \omterm{def:b2:expected-utility:ara}{coefficient of
absolute risk aversion} of traders with utility $-e^{-aW}$ over wealth $W$; $c$ the cost of a signal;
$\phi$ the share of \omterm{def:b2:prices-as-information:informed-trader}{informed traders}; $z$ the random net supply of the asset; $\xi$ the combination of
$\theta$ and $z$ that the price reveals; $\kappa$ the weight of $\xi$ in the forecast of $\theta$. $p$ is a
price and $\bar p$ a plan price; $\bar q$ a plan quota.
\end{notation}

\section{What a planner would need to know}\label{sec:b2:prices-as-information:planner}

The second welfare theorem says that any \omterm{def:b2:exchange-economies:pareto-efficient}{Pareto-efficient allocation} can be reached by markets after
transfers. A planner could also try to reach one directly, by telling every household what to consume
and every firm what to produce. To do it, the planner must find a point at which every household's
marginal rate of substitution between every pair of goods equals every firm's marginal rate of
transformation. That needs every household's preferences and every firm's technology.

% code: exercises/solutions.py counts
Count the numbers for an economy far simpler than any real one. Kestria has 2 million households and 40
goods, and every household has Cobb--Douglas preferences, so that one household is described by the 40
shares of its budget it spends on each good, which sum to one: 39 numbers. The planner must collect 78
million numbers before considering a single firm, and collect them again whenever tastes change. Worse,
households have no reason to report them truthfully if the report changes what they receive. A market
uses a different message: the prices of 39 goods in terms of the fortieth, which every household reads,
and each household's own purchases, which no one else needs to know.

\begin{definition}[Informational efficiency of an allocation mechanism]\label{def:b2:prices-as-information:informational-efficiency}
An \omterm{def:b2:exchange-economies:allocation}{allocation} mechanism has agents exchange messages until an \omterm{def:b2:exchange-economies:allocation}{allocation} is reached. It is
\emph{informationally efficient}\index{informational efficiency} if it reaches \omterm{def:b2:exchange-economies:pareto-efficient}{Pareto-efficient
allocations} with messages of the smallest dimension among all mechanisms that do so.
\end{definition}

\begin{omfigure}
\begin{tikzpicture}[econ, font=\footnotesize,
    box/.style={draw, rounded corners=2pt, minimum width=17mm, minimum height=6mm, align=center}]
  \node[box, fill=omDef!15] (pl) at (0,2.4) {Planner};
  \foreach \i/\x in {1/-1.9,2/0,3/1.9} {
    \node[box] (h\i) at (\x,0) {Agent \i};
    \draw[->, omThm, thick] (h\i.north) to[bend left=12] (pl.south);
    \draw[->, black!70] (pl.south) to[bend left=12] (h\i.north);
  }
  \node[align=center, omThm] at (-2.5,1.45) {preferences,\\technologies};
  \node[align=center] at (2.5,1.45) {quantities};
  \node[font=\small\bfseries] at (0,3.3) {A planner};
  \begin{scope}[xshift=7cm]
    \node[box, fill=omDef!15] (mk) at (0,2.4) {Prices $p$};
    \foreach \i/\x in {1/-1.9,2/0,3/1.9} {
      \node[box] (m\i) at (\x,0) {Agent \i};
      \draw[->, omThm, thick] (m\i.north) to[bend left=12] (mk.south);
      \draw[->, black!70] (mk.south) to[bend left=12] (m\i.north);
    }
    \node[align=center, omThm] at (-2.5,1.45) {own excess\\demand};
    \node[align=center] at (2.5,1.45) {prices};
    \node[font=\small\bfseries] at (0,3.3) {A market};
  \end{scope}
\end{tikzpicture}

{\small What travels in each mechanism. A planner must learn every agent's preferences and technology
and send back quantities. In a market each agent keeps its preferences to itself; it reads the prices
and answers with what it wants to buy or sell at them, and prices move until the answers add up to
zero.}
\end{omfigure}

\section{Prices as a short message}\label{sec:b2:prices-as-information:message}

In the competitive mechanism each consumer needs only two things to choose: her own preferences and the
price vector. Each firm needs its own technology and the prices. No agent needs anyone else's
information, and the prices, $n - 1$ numbers in an economy with $n$ goods, are the whole of what they share.
Whether some other mechanism could do with less is a mathematical question, and it has an answer.

\begin{theorem}[Minimal messages]\label{thm:b2:prices-as-information:minimal}
In exchange economies with smooth preferences, among the mechanisms that always reach \omterm{def:b2:exchange-economies:allocation}{allocations} that
are \omterm{def:b2:exchange-economies:pareto-efficient}{Pareto-efficient} and leave every consumer at least as well off as her \omterm{def:b2:exchange-economies:endowment}{endowment}, the competitive
mechanism uses messages of the smallest dimension, and it is the only one of that dimension (Jordan
1982).
\end{theorem}

\admitted

\begin{remark}
The claim that prices carry, in a few numbers, the information that an economy needs to coordinate
production and consumption is often called \emph{Hayek's result}. \Cref{thm:b2:prices-as-information:minimal}
is one precise version of it.
\end{remark}

The theorem is a statement about the static problem: given preferences and technologies, find an
efficient \omterm{def:b2:exchange-economies:allocation}{allocation}. It does not say that the prices are found quickly, or at all: the t\^atonnement of
\cref{ch:b2:exchange-economies} need not converge, and real markets find their prices through trading
institutions whose working \cref{ch:b2:exchange-economies} observed in the laboratory. Nor does it cover
information about the future, which is the subject of the rest of the chapter. A harvest, a company's
profits or an election result are unknown to everyone, and each trader holds a fragment of what is known
about them.

\section{Equilibrium when traders learn from prices}\label{sec:b2:prices-as-information:learning}

An asset will pay $v = s_1 + s_2 + \varepsilon$, where $s_1$, $s_2$ and $\varepsilon$ are independent and
normally distributed with mean 0 and variance 1. Half the traders observe $s_1$, the other half $s_2$; no
one observes $\varepsilon$. Every trader has utility $-e^{-aW}$ and the asset is in zero net supply. A
trader who believes $v$ has mean $m$ and variance $\sigma^2$ demands $x = (m - p)/(a\sigma^2)$ units
(\cref{exo:b2:prices-as-information:10} recalls why).

If traders ignore what the price says, a trader of the first group expects $v$ to be $s_1$ with variance
2, and demands $(s_1 - p)/(2a)$; the second group demands $(s_2 - p)/(2a)$. Equal groups clear the market
when the two demands sum to zero:
\begin{equation}\label{eq:b2:prices-as-information:naive}
p = \tfrac12\,(s_1 + s_2).
\end{equation}
The price depends on both signals, but it is not the best forecast of $v$, which is $s_1 + s_2$. And a
trader who knows the rule \eqref{eq:b2:prices-as-information:naive} and sees $p$ can compute $s_1 + s_2 =
2p$: the price tells her the other group's signal. A trader who ignores that is leaving information on the
table.

\begin{definition}[Fully revealing equilibrium]\label{def:b2:prices-as-information:fully-revealing}
A rational expectations equilibrium is a price function, from traders' information to prices,
such that markets clear when every trader chooses using her own information and what the price
function lets her infer from the price. It is \emph{fully revealing}\index{fully revealing equilibrium}
if the price tells every trader everything that all traders together know about the asset's value.
\end{definition}

\begin{proposition}[A fully revealing price]\label{prop:b2:prices-as-information:revealing}
In the two-signal market, $p = s_1 + s_2$ is a rational expectations equilibrium, and it is \omterm{def:b2:prices-as-information:fully-revealing}{fully
revealing}.
\end{proposition}

\begin{proof}
Suppose traders know the price function $p = s_1 + s_2$. A trader of the first group sees $p$ and $s_1$
and so knows $s_1 + s_2$; so does a trader of the second group. Both then expect $v$ to be $s_1 + s_2$, with
variance 1, and demand $(s_1 + s_2 - p)/a = 0$. Demand sums to zero, so the market clears. The price
reveals $s_1 + s_2$, which is all the traders together know.
\end{proof}

For example, if $s_1 = 1.2$ and $s_2 = 0.4$, the naive price is 0.8 and the \omterm{def:b2:prices-as-information:fully-revealing}{fully revealing} price is 1.6.

\begin{omfigure}
\begin{tikzpicture}[econ]
\begin{axis}[omrecord, grid=none, xmin=-4, xmax=4, ymin=-4, ymax=4, axis lines=middle,
    xlabel={$s_1 + s_2$}, ylabel={$p$}, xtick={-4,-2,2,4}, ytick={-4,-2,2,4},
    width=0.5\linewidth, height=0.5\linewidth, xlabel style={anchor=north west},
    ylabel style={anchor=south east}]
  \addplot[omDef, very thick, domain=-4:4, samples=2] {x};
  \addplot[omThm, thick, dashed, domain=-4:4, samples=2] {x/2};
  \node[anchor=west, font=\footnotesize, omDef] at (axis cs:1.4,3.4) {fully revealing};
  \node[anchor=north west, font=\footnotesize, omThm] at (axis cs:1.6,0.6) {naive};
  \node[eqpoint] at (axis cs:1.6,1.6) {};
  \node[eqpoint] at (axis cs:1.6,0.8) {};
\end{axis}
\end{tikzpicture}

{\small The price as a function of the aggregate signal. Both lines are increasing, so either price
would tell a trader who knows the line what $s_1 + s_2$ is. Only on the solid line does that knowledge leave
demands unchanged; on the dashed line, traders who learn from the price want to trade again. The dots
mark $s_1 + s_2 = 1.6$.}
\end{omfigure}

The proposition is a strong version of the claim that prices aggregate dispersed information: no one
knows $s_1 + s_2$, and the price shows it to everyone. It also contains a puzzle. In the \omterm{def:b2:prices-as-information:fully-revealing}{fully revealing
equilibrium} every trader's demand is zero whatever her own signal; the price contains the signals, yet no
trade depends on them. How the signals got into the price is not explained by the equilibrium, and if
signals had to be paid for, no one would pay.

\section{The limit: costly information}\label{sec:b2:prices-as-information:costly}

\begin{definition}[Informed trader]\label{def:b2:prices-as-information:informed-trader}
In a market with costly information, an \emph{informed trader}\index{informed trader} is one who has
paid the cost $c$ to observe a signal about the asset's value; an uninformed trader observes only the
price.
\end{definition}

Let the asset pay $v = \theta + \varepsilon$, with $\theta$ and $\varepsilon$ independent standard normal.
A share $\phi$ of the traders pays $c$ to observe $\theta$. The net supply of the asset is $z$, normal with
mean 0 and standard deviation $\sigma_z$, a stand-in for trades made for reasons unrelated to $v$
(a trader who needs cash, a fund that rebalances). Take $a = 1$. \omterm{def:b2:prices-as-information:informed-trader}{Informed traders} demand $\theta - p$; the
market clears when
\[
\phi(\theta - p) + (1 - \phi)\,x_U = z ,
\]
where $x_U$ is the uninformed demand, which depends on the price only. So the price depends on $\theta$ and
$z$ only through $\xi = \theta - z/\phi$, and an uninformed trader who sees $p$ learns $\xi$: a signal of
$\theta$ blurred by the noise $z/\phi$, whose standard deviation is $\sigma_z/\phi$. The more traders are
informed, the less the noise blurs the price. Her forecast of $\theta$ is $\kappa\xi$, with $\kappa = 1/(1 +
\sigma_z^2/\phi^2)$, and its error variance is $1 - \kappa$.

Who buys the signal? An \omterm{def:b2:prices-as-information:informed-trader}{informed trader} pays $c$ and forecasts $v$ with error variance 1; an uninformed
trader pays nothing and forecasts with error variance $\mathrm{Var}(v \mid p) = 1 + (1 - \kappa)$. With CARA
utility and normal returns, the ratio of their expected utilities works out to $e^{ac}\sqrt{\mathrm{Var}(v
\mid \theta)/\mathrm{Var}(v \mid p)}$ (Grossman and Stiglitz 1980). Traders are indifferent when this ratio
is one.

\begin{proposition}[The equilibrium share of informed traders]\label{prop:b2:prices-as-information:share}
With $a = 1$ and unit variances, let $k = e^{2c} - 1$ with $0 < k < 1$. Traders are indifferent between
buying the signal and not when the price forecasts $\theta$ with error variance
\[
\mathrm{Var}(\theta \mid p) = k ,
\]
which happens at the informed share $\phi^* = \sigma_z\sqrt{(1 - k)/k}$ when this is below one. If it is
above one, every trader buys the signal.
\end{proposition}

\begin{proof}
Indifference requires $e^{2c}\cdot 1 = \mathrm{Var}(v \mid p) = 1 + \mathrm{Var}(\theta \mid p)$, that is,
$\mathrm{Var}(\theta \mid p) = e^{2c} - 1 = k$. The price reveals $\xi = \theta + u$ with $u$ of variance
$\sigma_u^2 = \sigma_z^2/\phi^2$, so $\mathrm{Var}(\theta \mid p) = \sigma_u^2/(1 + \sigma_u^2)$. Setting this to
$k$ gives $\sigma_u^2 = k/(1 - k)$ and $\phi = \sigma_z/\sigma_u = \sigma_z\sqrt{(1 - k)/k}$. If more traders
were informed, the price would be more revealing, the uninformed would do better than the informed, and
some would stop buying; if fewer, the reverse.
\end{proof}

% code: exercises/solutions.py gs
With $\sigma_z = 0.25$ and $c = 0.1116$, $k = e^{0.2232} - 1 = 0.250$, $\sigma_u = \sqrt{0.25/0.75} = 0.577$
and $\phi^* = 0.25/0.577 = 0.433$: 43.3 per cent of the traders buy the signal, and the price forecasts
$\theta$ with error variance 0.25, a quarter of the variance of $\theta$ itself. The listing checks the
result by brute force. It draws 400{,}000 periods of the market at that informed share, computes the
price from the clearing condition, and compares the realised utilities of the two kinds of trader.

\omcode{code/bachelor-2/08-prices-as-information/listings/informed_share.py}{1}{30}{A market with
informed and uninformed traders, simulated at the equilibrium informed
share.}\label{lst:b2:prices-as-information:informed-share}

\omoutput{code/bachelor-2/08-prices-as-information/outputs/informed_share.txt}{The price forecasts
$\theta$ with error variance 0.25; the two kinds of trader are almost equally well
off.}\label{out:b2:prices-as-information:informed-share}

% code: listings/informed_share.py
The two utilities, $-0.9917$ and $-0.9922$, differ by less than the sampling noise of the simulation. The
fitted slope 0.862 is the uninformed trader's reading of the price, $\kappa/b$ in the listing's notation.

\begin{omfigure}
\begin{tikzpicture}[econ]
\begin{axis}[omrecord, grid=none, xmin=0, xmax=0.4, ymin=0, ymax=1.05, xlabel={Cost of the signal $c$},
    ylabel={Share}, width=0.62\linewidth, height=0.42\linewidth, xtick={0,0.1,0.2,0.3,0.4},
    xticklabel style={/pgf/number format/fixed}, legend columns=2,
    legend style={at={(0.5,-0.25)}, anchor=north, font=\footnotesize}]
  \addplot[omDef, very thick] table[x=c, y=phi]
    {figdata/bachelor-2/08-prices-as-information/out/gs.csv};
  \addplot[omThm, thick, dashed] table[x=c, y=resid]
    {figdata/bachelor-2/08-prices-as-information/out/gs.csv};
  \legend{Informed share $\phi^*$, Error variance $\mathrm{Var}(\theta \mid p)$}
\end{axis}
\end{tikzpicture}

{\small The Grossman--Stiglitz market with $\sigma_z = 0.25$. As the signal becomes dearer, fewer traders
buy it and the price forecasts $\theta$ less well. Below $c = 0.029$ every trader is informed; above $c =
0.347$ no one is, and the price says nothing about $\theta$.}
\end{omfigure}

Two features of the proposition are worth stating. The price's error variance $k$ depends only on the cost
of information, not on the noise: more noise in supply attracts more \omterm{def:b2:prices-as-information:informed-trader}{informed traders}, exactly enough to
restore the same informativeness. And the price can never be \omterm{def:b2:prices-as-information:fully-revealing}{fully revealing} while information costs
something.

\begin{definition}[Grossman--Stiglitz paradox]\label{def:b2:prices-as-information:gs-paradox}
The \emph{Grossman--Stiglitz paradox}\index{Grossman--Stiglitz paradox} is the result that when
information is costly, prices cannot fully reveal it: if they did, no trader would pay for information,
and then prices would contain none.
\end{definition}

Without noise ($\sigma_z = 0$) there is no equilibrium at all (\cref{exo:b2:prices-as-information:11}). With
noise, the market settles where prices are informative enough to make information barely worth its cost.
Prices are then partly revealing, and the informed earn, on average, exactly enough to pay for what they
know.

\begin{model}{Grossman--Stiglitz}
\assumes Traders with CARA utility; normal values and noise; a fixed cost of information; traders who
learn from prices.
\predicts An equilibrium share of \omterm{def:b2:prices-as-information:informed-trader}{informed traders}; prices partly revealing, more revealing when
information is cheaper; returns to being informed that just cover the cost.
\tested No direct test is reported here: a test would need the informed share and the cost of information,
which are rarely observed. The model is consistent with the racetrack and election records of this chapter,
where prices carry much but not all of the information.
\breaks Normality and CARA are chosen for tractability; noise traders are assumed, not explained; the
informed share is hard to measure.
\end{model}

\section{Prediction markets and their record}\label{sec:b2:prices-as-information:prediction}

\begin{definition}[Prediction market]\label{def:b2:prices-as-information:prediction-market}
A \emph{prediction market}\index{prediction market} trades contracts whose payoff depends on a future
event, such as a contract that pays one dollar if a candidate wins, or one cent per percentage point of
her vote share. The price of such a contract is read as the market's forecast of the probability or the
quantity.
\end{definition}

A contract that pays one dollar if an event occurs is worth, to a risk-neutral trader, the probability of
the event. If traders hold scattered information about the event and the market aggregates it, as in
\cref{prop:b2:prices-as-information:revealing}, the price is a better forecast than any single trader's.
The racetrack of the opening chart is such a market, and its record shows one systematic error.

\begin{definition}[Favourite--longshot bias]\label{def:b2:prices-as-information:longshot-bias}
The \emph{favourite--longshot bias}\index{favourite--longshot bias} is the regularity that, in betting
markets, bets on outcomes with low probability (longshots) return less on average than bets on outcomes
with high probability (favourites): prices overstate the chance of unlikely outcomes and understate the
chance of likely ones.
\end{definition}

% code: exercises/solutions.py payback
Odds of 100/1 mean that a winning dollar bet returns 101 dollars. If such horses won one race in 101, the
expected return would be zero: the odds are already net of the track's share. Unbiased odds would return
what the track leaves, about $-18$ to $-23$ per cent, as for horses at 4/1 to 9/1 and for a horse picked at
random, so a 100/1 horse would win about $0.77/101 = 0.0076$ of its races. A rate of return of $-61$ per cent
says instead that such horses win about $0.39/101 = 0.0039$ of their races: about half as often as that
benchmark implies.

\begin{record}[Election markets against polls]
% ledger: b2:bfnr2008-polls
The Iowa Electronic Markets trade vote-share contracts on elections. Berg, Forsythe,
Nelson and Rietz (2008) compare 237 contract prices, in 49 markets in 13 countries, with the vote shares
that followed. In 15 national elections with polls to compare, the market's forecast was closer to the
result than the polls' average in 9 of the 15. Across those elections the average absolute error of the
final polls was 1.91 percentage points of the vote; the market's was 1.49 points with election-eve prices
and 1.58 points with prices averaged over the last week. Their markets did better for presidential
elections than for lower-profile ones, and with more trading volume near the election.
\end{record}

The record supports aggregation and qualifies it. The market's errors are a little smaller than the
polls', not dramatically so, and the markets studied are those that existed and drew traders. A market
can also be pushed: a trader with deep pockets can buy a contract to move its price. Whether such a push
lasts is itself a test: if other traders treat the moved price as a profit opportunity, it reverts
(\cref{exo:b2:prices-as-information:9}).

\begin{model}{Prices as information}
\assumes Traders hold dispersed, partial information; they trade on it and learn from prices; markets
are thick enough for prices to respond.
\predicts Prices that aggregate dispersed signals, forecasting better than any single participant;
unpredictable price changes once the information is in.
\tested \omterm{def:b2:prices-as-information:prediction-market}{Prediction markets} forecast elections about as well as or better than the polls' average
(Berg et al.\ 2008); racetrack odds rank horses' chances correctly but overstate those of longshots
(Snowberg and Wolfers 2010); laboratory markets (\cref{ch:b2:exchange-economies}).
\breaks Thin markets, manipulation, systematic misperception of small probabilities (the
\omterm{def:b2:prices-as-information:longshot-bias}{favourite--longshot bias}), and costly information (Grossman--Stiglitz).
\end{model}

\section{Planned prices as data}\label{sec:b2:prices-as-information:plan}

A planned economy also has prices: the prices at which state enterprises deliver goods to one another
and to shops. Set by an office rather than by trade, they need not clear markets. When a plan price is
below the price that would clear the market, the model of a price ceiling predicts that buyers want more
than sellers supply; the goods are rationed, by queue, by permit or by connection. The plan price then
carries the planner's information, not the buyers' and sellers'.

\begin{definition}[Dual-track pricing]\label{def:b2:prices-as-information:dual-track}
\emph{Dual-track pricing}\index{dual-track pricing} is a reform in which producers must still deliver
fixed plan quotas at plan prices, but may sell any output above the quota, and buy inputs beyond their
plan \omterm{def:b2:exchange-economies:allocation}{allocation}, at market prices.
\end{definition}

Take a mine with cost $C(q) = q^2/2$ that must deliver a quota $\bar q = 18$ at the plan price $\bar p = 20$,
and may sell anything more at the market price $p = 30$. Its profit is $\bar p\bar q + p(q - \bar q) - C(q)$, so
the last tonne is worth $p$ and it produces where $C'(q) = q = 30$. It delivers 18 tonnes to the plan and sells
12 on the market, for a profit of $360 + 360 - 450 = 270$. With full liberalisation it would produce the
same 30 tonnes and earn $900 - 450 = 450$. The difference, $(p - \bar p)\bar q = 180$, is a transfer from the
mine to the buyers entitled to plan deliveries, and it does not depend on anything the mine chooses.

\begin{proposition}[The plan track as a lump-sum transfer]\label{prop:b2:prices-as-information:dual-track}
Under \omterm{def:b2:prices-as-information:dual-track}{dual-track pricing}, if the quota is enforced and a producer's output exceeds it, the producer's
output choice is the one it would make at the market price alone, and the plan track transfers $(p - \bar
p)\bar q$ from the producer to the holders of plan rights, a \omterm{def:b2:welfare-theorems:lump-sum}{lump-sum transfer}.
\end{proposition}

\begin{proof}
Profit is $\bar p\bar q + p(q - \bar q) - C(q) = [pq - C(q)] - (p - \bar p)\bar q$. The bracket is profit at the
market price; the second term does not depend on $q$. So the $q$ that maximises profit is the one that
maximises profit at the market price, and the producer pays $(p - \bar p)\bar q$ whatever it does.
\end{proof}

This is the interpretation of Lau, Qian and Roland (2000): the plan track, kept and enforced, protects
those who held rights under the plan, while the market track puts every decision at the margin under
market prices; the reform can then be Pareto-improving, an application of the \omterm{def:b2:welfare-theorems:lump-sum}{lump-sum transfers} of
\cref{ch:b2:welfare-theorems}. The proposition also says what can go wrong: if quotas are not enforced,
or plan goods can be resold, the transfer is no longer lump-sum.

\begin{record}[Two prices for the same goods: China's dual track]
% ledger: b2:zhang2018-text
From January 1985, Chinese state enterprises could sell output above their quotas at market prices, while
plan supplies stayed at the 1983 quota. A survey of 429 state-owned industrial enterprises found that
the output they sold themselves rose from 30.1 per cent of their sales in 1984 to 43.8 per cent in the
first half of 1985. Coal output rose by 105.1 million tonnes between 1982 and 1984, and 53.2 per cent of
the increase came from rural and township mines, outside the central plan (Zhang 2018).
\end{record}

\begin{omchart}{Shares of China's coal output by type of mine, in per cent, 1978--1984.}{Zhang (2018),
  Table~12.2}
\begin{tikzpicture}
\begin{axis}[omrecord, xmin=1977.5, xmax=1984.5, ymin=0, ymax=60, xtick={1978,...,1984},
    xticklabel style={/pgf/number format/1000 sep={}}, ylabel={Per cent of national output},
    width=0.62\linewidth, height=0.42\linewidth, legend columns=3,
    legend style={at={(0.5,-0.2)}, anchor=north, font=\footnotesize}]
  \addplot[omDef, very thick, mark=*] table[x=year, y=central]
    {figdata/bachelor-2/08-prices-as-information/out/coal.csv};
  \addplot[black!60, thick, mark=square*] table[x=year, y=local]
    {figdata/bachelor-2/08-prices-as-information/out/coal.csv};
  \addplot[omThm, very thick, mark=triangle*] table[x=year, y=nonstate]
    {figdata/bachelor-2/08-prices-as-information/out/coal.csv};
  \legend{Centrally allocated, Local government, Non-state}
\end{axis}
\end{tikzpicture}
\end{omchart}

% figdata: coal_summary (central -7, local -4.1, nonstate +10.6)
Between 1978 and 1984 the non-state mines' share rose by 10.6 points, from 15.6 to 26.2 per cent, while the
centrally allocated mines' share fell by 7.0 points. In the model, a mine outside the plan responds only
to the market price, so its growth is what the market price called for. The chart cannot separate that
response from other reforms of the same years, and the shares are of output, not of welfare. It shows
the plan's prices and the market's at work side by side, which is what a single price system cannot show.

Planned prices fail as messages for the reason of \cref{sec:b2:prices-as-information:planner}: the office
setting them lacks the information that buyers and sellers hold, and those who hold it have no reason to
report it. Market prices fail in the ways of this chapter: when information is costly, when markets are
thin, when traders misjudge small probabilities. Neither failure is evidence that the other system works.

\section{Exercises}\label{sec:b2:prices-as-information:exercises}

\begin{exercise}[$\star$]\label{exo:b2:prices-as-information:1}
Read the opening chart. How many cents does a dollar bet on the favourite return on average, and a dollar
bet at odds of 100/1 or longer? If odds were unbiased forecasts, what would the chart look like?
\end{exercise}

\begin{exercise}[$\star$]\label{exo:b2:prices-as-information:2}
An economy has 2 million households and 40 goods, and each household has Cobb--Douglas preferences. How
many numbers must a planner collect to know all the preferences? How many prices does a market post?
What else would the planner need?
\end{exercise}

\begin{exercise}[$\star$]\label{exo:b2:prices-as-information:3}
In the two-signal market, $s_1 = 2$ and $s_2 = -0.5$. Find the naive price and the \omterm{def:b2:prices-as-information:fully-revealing}{fully revealing} price.
At the naive price, would traders who learnt $s_1 + s_2$ from the price want to buy or to sell?
\end{exercise}

\begin{exercise}[$\star$]\label{exo:b2:prices-as-information:4}
Describe data from a \omterm{def:b2:prices-as-information:prediction-market}{prediction market} that would refute the claim that its prices aggregate the
information of its traders.
\end{exercise}

\begin{exercise}[$\star\star$]\label{exo:b2:prices-as-information:5}
In the Grossman--Stiglitz market with $\sigma_z = 0.25$, find the equilibrium informed share and the
error variance of the price's forecast when $c = 0.05$.
\end{exercise}

\begin{exercise}[$\star\star$]\label{exo:b2:prices-as-information:6}
Predict, then compute: the cost of the signal rises from 0.1116 to 0.2. What happens to the informed share
and to the informativeness of the price? What would happen if instead $\sigma_z$ doubled?
\end{exercise}

\begin{exercise}[$\star\star$]\label{exo:b2:prices-as-information:7}
Use the record box on election markets. By what percentage was the market's average error smaller than the
polls', by each measure? In what share of elections was the market closer? What would you need to know to
judge whether 9 of 15 could be luck?
\end{exercise}

\begin{exercise}[$\star\star$]\label{exo:b2:prices-as-information:8}
A favourite at odds of 2/1 returns 3 dollars per dollar bet when it wins. If favourites return $-5.5$ per
cent on average, how often do such horses win? What probability do the odds suggest? Compare with the
horse at 100/1 in the text.
\end{exercise}

\begin{exercise}[$\star\star$]\label{exo:b2:prices-as-information:9}
A trader buys a large quantity of a candidate's contract in a \omterm{def:b2:prices-as-information:prediction-market}{prediction market}, raising its price from
0.40 to 0.55 within an hour, with no news about the election. Design a test of whether the market's price
was manipulated, and say what each outcome would show.
\end{exercise}

\begin{exercise}[$\star\star\star$]\label{exo:b2:prices-as-information:10}
(a) A trader with utility $-e^{-aW}$ buys $x$ units of an asset at price $p$ whose value is normal with mean
$m$ and variance $\sigma^2$. Using $\mathbb{E}[e^{-Y}] = e^{-\mathbb{E}Y + \mathrm{Var}(Y)/2}$ for normal $Y$, show
that her best $x$ is $(m - p)/(a\sigma^2)$. (b) In the two-signal market with signal variances $\sigma_1^2$ and
$\sigma_2^2$ and the noise variance $\sigma_\varepsilon^2$, show that $p = s_1 + s_2$ is still a \omterm{def:b2:prices-as-information:fully-revealing}{fully
revealing equilibrium}.
\end{exercise}

\begin{exercise}[$\star\star\star$]\label{exo:b2:prices-as-information:11}
Show that in the Grossman--Stiglitz market without noise ($\sigma_z = 0$) and with $0 < c < \tfrac12\ln 2$
there is no equilibrium: neither $\phi = 0$ nor any $\phi > 0$ leaves traders indifferent or content with
their choice.
\end{exercise}

\begin{exercise}[$\star\star\star$]\label{exo:b2:prices-as-information:12}
In the equilibrium of the text ($\sigma_z = 0.25$, $c = 0.1116$), how much do traders spend on
information, per trader? The asset's payoff does not depend on who is informed. Who gains from the
information, and is spending on it a social gain or a transfer? Under what circumstances would better
prices have real value?
\end{exercise}

\section{Problem: Kestria's harvest-forecast market}\label{sec:b2:prices-as-information:problem}

\begin{problem}[{Weekend problem --- Kestria, an invented economy, runs a market in contracts on next
year's wheat harvest: how well does the price forecast the harvest, how many traders pay for a forecast,
and does the market beat the national forecast office}]
\label{pb:b2:prices-as-information:1}
Kestria's wheat harvest, in thousands of tonnes, will be $H = 100 + s_N + s_S + \varepsilon$. Farmers in the
north see $s_N$, farmers in the south see $s_S$; the two signals have mean 0 and variance 16, $\varepsilon$
has mean 0 and variance 9, and all three are independent and normal.

\medskip
\noindent\textbf{Part I --- Signals and a planner.}
\begin{enumerate}
  \item Find the variance of $H$.
  \item A planner who learns both signals forecasts $H$ with which error variance? With one signal? With
        none?
  \item How many numbers must the planner collect each year, and from whom?
  \item Why might northern farmers report a low $s_N$ if the planner sets the price of seed by the
        forecast?
  \item A market forecast that used both signals would have which error variance?
  \item Which result of the chapter says a market price could achieve it?
\end{enumerate}

\medskip
\noindent\textbf{Part II --- The market price.} A contract pays $H$ at harvest. Northern and southern
traders, equal in number, have utility $-e^{-W}$; the contract is in zero net supply. This year $s_N = 6$ and
$s_S = -2$.
\begin{enumerate}[resume]
  \item If traders ignore the price, find each group's demand and the price.
  \item At that price, how much does a northern trader buy?
  \item If traders learn from the price, find the equilibrium price.
  \item Check that the market clears at that price.
  \item Which price is the better forecast of $H$, and what is its error variance?
  \item What does a southern trader learn from the price of question 9?
\end{enumerate}

\medskip
\noindent\textbf{Part III --- Costly information.} Use the Grossman--Stiglitz market of
\cref{sec:b2:prices-as-information:costly}, with $\sigma_z = 0.25$: a forecast of the harvest costs $c =
0.08$.
\begin{enumerate}[resume]
  \item Compute $k$.
  \item Find the equilibrium informed share.
  \item Find the error variance of the price's forecast of $\theta$.
  \item The cost falls to 0.05. What changes?
  \item Explain why the price cannot reveal $\theta$ fully.
  \item Who pays for the information the uninformed traders read from the price?
\end{enumerate}

\medskip
\noindent\textbf{Part IV --- The market against the forecast office.} The office's forecast is $100 + s_N +
s_S$ plus a survey error of variance 4; the market's price is $100 + s_N + s_S$ plus noise of variance 2.25.
The table records twelve seasons.

\begin{center}
\small
\begin{tabular}{lrrrrrrrrrrrr}
\toprule
Season & 1 & 2 & 3 & 4 & 5 & 6 & 7 & 8 & 9 & 10 & 11 & 12 \\
\midrule
Harvest & 93.2 & 118.4 & 108.6 & 91.8 & 101.5 & 109.7 & 86.3 & 106.8 & 96.5 & 95.9 & 100.1 & 96.4 \\
Office  & 90.9 & 110.2 & 101.2 & 95.5 & 102.9 & 109.9 & 84.4 & 103.6 & 99.4 & 98.0 & 100.4 & 99.6 \\
Market  & 94.8 & 113.4 & 101.9 & 95.9 & 102.2 & 103.8 & 86.7 & 105.5 & 98.2 & 92.2 & 103.5 & 100.5 \\
\bottomrule
\end{tabular}
\end{center}

\begin{enumerate}[resume]
  \item Find the expected squared error of each forecast.
  \item Compute the mean absolute error of each forecast over the twelve seasons, and the mean squared
        error.
  \item In how many seasons was the market closer?
  \item The difference between the market's and the office's squared errors has mean $-0.72$ and standard
        deviation 17.8 over the twelve seasons. Is the record informative?
  \item Compute the mean absolute error of the average of the two forecasts, and its expected squared
        error.
  \item State the result: the informed share and the price's error variance in Part III, and what the
        twelve seasons show.
\end{enumerate}
\end{problem}
